\chapter{Pourquoi il existe une fourchette}\label{ch:mx:why-there-is-a-spread}

Un teneur de marché affiche un cours acheteur de 100.00 et un cours vendeur de 100.02. Si personne n'en savait plus qu'elle, les deux centimes rémunéreraient son temps, ses systèmes et le
risque de l'inventaire qu'elle conserve entre deux transactions. Un acheteur qui a lu l'actualité avant les autres fait servir la fourchette à autre chose~: chaque transaction
avec lui fait perdre de l'argent, et la fourchette doit récupérer cette perte auprès de tous les autres. Trois théories de la fourchette répondent à trois questions~: combien
coûte le traitement d'une transaction, combien coûte la détention de ce que la transaction laisse derrière elle, et combien coûte la négociation avec quelqu'un qui en sait davantage. Elles conduisent
à des fourchettes différentes, à des dynamiques de cotations différentes et à des manières différentes de mesurer chaque composante (chapitre~5).

\section{Coûts de traitement des ordres}

\begin{definition}[Coût de traitement des ordres]\label{def:mx:why-there-is-a-spread:processing}
Le \emph{coût de traitement des ordres}\index{coût de traitement des ordres} d'un fournisseur de liquidité est le coût par transaction qui ne dépend ni de
l'identité de la contrepartie ni de ce qui arrive ensuite au cours~: frais de bourse et de compensation, systèmes, personnel et capital, nets des remises.
\end{definition}

Si le traitement était le seul coût, un teneur de marché en concurrence afficherait une demi-fourchette égale à ce coût, la valeur ne réagirait pas aux transactions, et
les cours de transaction oscilleraient entre le cours acheteur et le cours vendeur autour d'une valeur suivant une marche aléatoire. Roll (1984) a transformé cette oscillation en estimation~: avec une
demi-fourchette effective constante $c$ et des signes de la transaction indépendants, les variations de prix ont une autocovariance de $-c^2$, donc $s = 2\sqrt{-\mathrm{Cov}(\Delta p_t,\Delta
p_{t-1})}$ (estimateur de Roll, livre~4, chapitre~21). Sur l'heure simulée de \texttt{firm.tape}, l'estimation vaut 0.88 ticks contre une fourchette affichée
(quoted spread) de 1.08~: les signes de la transaction du simulateur sont positivement autocorrélés (0.29 au retard un, car les ordres de bruit arrivent en métaordres), ce qui retire
une partie de l'oscillation, comme c'est le cas sur les marchés réels (livre~7, chapitre~9).

\section{Inventaire}

\begin{definition}[Coût de détention de l'inventaire]\label{def:mx:why-there-is-a-spread:inventory}
Le \emph{coût de détention de l'inventaire}\index{coût de détention de l'inventaire} d'une transaction est la compensation qu'un fournisseur de liquidité averse au risque exige pour porter
la position que la transaction lui laisse jusqu'à pouvoir la déboucler.
\end{definition}

Garman (1976) a modélisé un dealer dont l'inventaire suit une marche aléatoire, pilotée par l'arrivée d'acheteurs et de vendeurs, et qui doit fixer ses prix pour éviter la ruine~;
Stoll (1978) et Ho et Stoll (1981) ont valorisé le risque. La version la plus simple contient déjà les deux caractéristiques que partagent tous les modèles d'inventaire.

\begin{proposition}[Dealer moyenne--variance]\label{prop:mx:why-there-is-a-spread:inventory}
Un dealer valorise son inventaire $q$ à la valeur $v$, s'attend à le conserver pendant un temps $T$ durant lequel la valeur a pour variance $\sigma^2T$, et a des préférences moyenne--variance
avec une aversion au risque $\gamma$~: détenir $q$ lui coûte $\tfrac12\gamma\sigma^2Tq^2$. Les prix auxquels une transaction de taille $Q$ le laisse indifférent
sont
\[
a = v + \gamma\sigma^2T\left(\frac Q2 - q\right),\qquad b = v - \gamma\sigma^2T\left(\frac Q2 + q\right):
\]
une fourchette $\gamma\sigma^2TQ$ qui croît avec la taille de la transaction, centrée sur le prix de réserve $v-\gamma\sigma^2Tq$ qui se déplace à l'encontre de son inventaire.
\end{proposition}

\begin{proof}
Vendre $Q$ au prix $a$ modifie sa richesse de $Q(a-v)$ et fait passer son coût de détention de $\tfrac12\gamma\sigma^2Tq^2$ à $\tfrac12\gamma\sigma^2T(q-Q)^2$~;
l'indifférence donne $Q(a-v)=\tfrac12\gamma\sigma^2T\left((q-Q)^2-q^2\right)=\gamma\sigma^2T\left(\tfrac{Q^2}{2}-qQ\right)$. Le cours acheteur s'en déduit avec $-Q$.
\end{proof}

Un dealer qui est long abaisse les deux cotations~: il veut vendre et décourage les achats des autres. Le terme d'inventaire est à la racine des modèles de teneur de marché
du livre~11 (le prix de réserve d'Avellaneda--Stoikov a la même forme)~; ce qui importe ici est qu'il produise une fourchette sans la moindre
information.

\section{Information~: le modèle à transactions séquentielles}

\begin{definition}[Trader informé, noise trader]\label{def:mx:why-there-is-a-spread:informed}
Un \emph{trader informé}\index{trader informé} connaît quelque chose sur la valeur que le fournisseur de liquidité ignore et trade en fonction de cette information. Un \emph{noise
trader}\index{noise trader} (trader de liquidité) trade pour des raisons sans rapport avec la valeur (un flux de fonds, une couverture, un rééquilibrage) et ne porte aucune information.
\end{definition}

\begin{definition}[Modèle à transactions séquentielles, modèle de Glosten--Milgrom]\label{def:mx:why-there-is-a-spread:gm}
Un \emph{modèle à transactions séquentielles}\index{modèle à transactions séquentielles} fait arriver les traders un par un auprès d'un fournisseur de liquidité neutre au risque et en concurrence, qui
affiche ses cotations avant chaque arrivée et apprend de chaque transaction. Le \emph{modèle de Glosten--Milgrom}\index{modèle de Glosten--Milgrom} est le modèle à transactions séquentielles dans lequel
chaque arrivée est informée avec la probabilité $\mu$ et trade dans le sens de la valeur, ou bien est un \omterm{def:mx:why-there-is-a-spread:informed}{noise trader} qui achète ou vend avec la même probabilité, et
le fournisseur fixe le cours vendeur et le cours acheteur à l'espérance de la valeur conditionnelle à un achat et à une vente.
\end{definition}

\begin{omfigure}
\begin{tikzpicture}[font=\footnotesize,grow=right,level distance=2.6cm,
  level 1/.style={sibling distance=2.6cm},level 2/.style={sibling distance=1.25cm},
  every node/.style={align=center}]
\node[draw,rounded corners] {valeur\\ $v_H$ ou $v_L$}
  child {node[draw,rounded corners,fill=omThm!12] {bruit\\ $1-\mu$}
    child {node {vend~: $\tfrac12$}}
    child {node {achète~: $\tfrac12$}}
    edge from parent node[below,pos=0.55] {}}
  child {node[draw,rounded corners,fill=omDef!12] {informé\\ $\mu$}
    child {node {vend si $v=v_L$}}
    child {node {achète si $v=v_H$}}};
\end{tikzpicture}

\omcaption{Une arrivée dans le \omterm{def:mx:why-there-is-a-spread:gm}{modèle de Glosten--Milgrom}. Le fournisseur ne voit que le sens de la transaction~; un achat est plus probable lorsque la valeur est élevée,
de $\mu$, et le cours vendeur est l'espérance de la valeur conditionnelle à un achat.}\label{fig:mx:why-there-is-a-spread:tree}
\end{omfigure}

\begin{proposition}[Cotations de Glosten--Milgrom]\label{prop:mx:why-there-is-a-spread:gmquotes}
Avec $v\in\{v_L,v_H\}$, $\Pr(v=v_H)=\delta$ et une part de \omterm{def:mx:why-there-is-a-spread:informed}{traders informés} $\mu$,
\[
a = v_L+(v_H-v_L)\frac{\delta(1+\mu)}{\delta(1+\mu)+(1-\delta)(1-\mu)},\qquad
b = v_L+(v_H-v_L)\frac{\delta(1-\mu)}{\delta(1-\mu)+(1-\delta)(1+\mu)}.
\]
Pour $\delta=\tfrac12$, la fourchette vaut $a-b=\mu(v_H-v_L)$. Après un achat, la croyance du fournisseur devient $\delta'=a$ (sur l'échelle $v_L=0$, $v_H=1$), et les
cotations sont sans regret~: elles correspondent à ce que chaque transaction révèle.
\end{proposition}

\begin{proof}
$\Pr(\text{achat}\mid v_H)=\mu+\tfrac{1-\mu}{2}=\tfrac{1+\mu}{2}$ et $\Pr(\text{achat}\mid v_L)=\tfrac{1-\mu}{2}$. La règle de Bayes (livre~2, chapitre~30,
pour la mise à jour bayésienne) donne $\Pr(v_H\mid\text{achat})$, et $a=v_L+(v_H-v_L)\Pr(v_H\mid\text{achat})$ par neutralité au risque et profit espéré nul. Le cours acheteur est
symétrique. Pour $\delta=\tfrac12$, $a=v_L+(v_H-v_L)\tfrac{1+\mu}{2}$ et $b=v_L+(v_H-v_L)\tfrac{1-\mu}{2}$.
\end{proof}

\omcode{code/firm/spreadmodels/firm_spreadmodels.py}{29}{39}{Cotations de Glosten--Milgrom et mise à jour après une transaction.}\label{lst:mx:why-there-is-a-spread:gm}

La fourchette relève entièrement de la sélection adverse (livre~1, chapitre~1)~: le fournisseur équilibre ses comptes entre les traders, perdant face aux informés exactement ce qu'il
gagne avec les \omterm{def:mx:why-there-is-a-spread:informed}{noise traders}. Elle se resserre à mesure que les transactions révèlent la valeur (\cref{fig:mx:why-there-is-a-spread:gm})~: avec 30~\% d'arrivées informées, la trajectoire médiane met 19
transactions pour que la fourchette tombe au dixième de son point de départ, avec 10~\%, elle en met 148. Le trading informé fait deux choses à la fois~: il élargit la fourchette et accélère
l'apprentissage des prix. Avec un coût de traitement $c$ par transaction ajouté à chaque cotation, la part de la fourchette due à la sélection adverse vaut $\mu\Delta v/(\mu\Delta
v+2c)$~: avec $\Delta v=1$ et $c=0.05$, la moitié de la fourchette pour $\mu=0.1$ et les trois quarts pour $\mu=0.3$.

\begin{omfigure}
\begin{tikzpicture}
\begin{axis}[width=9cm,height=5.2cm,xlabel={\small arrivées},ylabel={\small fourchette moyenne ($v_H-v_L=1$)},tick label style={font=\footnotesize},
  xmin=0,xmax=150,ymin=0,ymax=0.52,grid=major,grid style={gray!25},legend columns=3,legend style={font=\footnotesize,at={(0.5,-0.3)},anchor=north,
  draw=none,fill=none,/tikz/every even column/.append style={column sep=8pt}},table/col sep=comma]
\addplot[omDef,very thick] table[x=n,y=mu10]{figdata/microstructure/04-why-there-is-a-spread/gm.csv};
\addplot[omProp,very thick] table[x=n,y=mu30]{figdata/microstructure/04-why-there-is-a-spread/gm.csv};
\addplot[omThm,very thick] table[x=n,y=mu50]{figdata/microstructure/04-why-there-is-a-spread/gm.csv};
\legend{$\mu=0.1$,$\mu=0.3$,$\mu=0.5$}
\end{axis}
\end{tikzpicture}

\omcaption{La fourchette de Glosten--Milgrom le long d'une suite d'arrivées, moyennée sur 400 trajectoires (la moitié avec une valeur haute)~: elle part de $\mu$ et décroît
à mesure que les transactions révèlent la valeur, plus vite lorsque davantage d'arrivées sont informées. Données~: \texttt{mx\_spread.gm\_learning}.}\label{fig:mx:why-there-is-a-spread:gm}
\end{omfigure}

\section{Information~: le trader stratégique}

\begin{definition}[Modèle de Kyle]\label{def:mx:why-there-is-a-spread:kyle}
Le \emph{modèle de Kyle}\index{modèle de Kyle} (Kyle model) comporte un \omterm{def:mx:why-there-is-a-spread:informed}{trader informé} qui connaît la valeur $v\sim\mathcal N(p_0,\Sigma_0)$ et choisit une quantité $x$, des \omterm{def:mx:why-there-is-a-spread:informed}{noise
traders} qui soumettent $u\sim\mathcal N(0,\sigma_u^2)$ indépendant de $v$, et des teneurs de marché concurrentiels et neutres au risque qui ne voient que le flux d'ordres total
$y=x+u$ et fixent le prix $p=\E[v\mid y]$.
\end{definition}

Le \omterm{def:mx:why-there-is-a-spread:informed}{trader informé} est stratégique~: il sait que ses transactions déplacent le prix et se cache dans le bruit. L'équilibre linéaire se déduit sans peine.

\begin{proposition}[Équilibre à une période de Kyle]\label{prop:mx:why-there-is-a-spread:kyle}
Dans l'unique équilibre linéaire, $x=\beta(v-p_0)$ et $p=p_0+\lambda y$ avec
\[
\beta=\frac{\sigma_u}{\sqrt{\Sigma_0}},\qquad \lambda=\frac{\sqrt{\Sigma_0}}{2\sigma_u};
\]
le prix révèle la moitié de l'information privée ($\mathrm{Var}(v\mid p)=\Sigma_0/2$) et le \omterm{def:mx:why-there-is-a-spread:informed}{trader informé} gagne
$\tfrac12\sqrt{\Sigma_0}\,\sigma_u$ en espérance, payé par les \omterm{def:mx:why-there-is-a-spread:informed}{noise traders}.
\end{proposition}

\begin{proof}
Étant donné $p=p_0+\lambda y$, le \omterm{def:mx:why-there-is-a-spread:informed}{trader informé} maximise $\E[x(v-p)\mid v]=x(v-p_0)-\lambda x^2$, d'où $x=(v-p_0)/(2\lambda)$ et $\beta=1/(2\lambda)$. Étant donné
$x=\beta(v-p_0)$, la projection du teneur de marché donne $\lambda=\mathrm{Cov}(v,y)/\mathrm{Var}(y)=\beta\Sigma_0/(\beta^2\Sigma_0+\sigma_u^2)$. En substituant
$\beta=1/(2\lambda)$, on obtient $\beta^2\Sigma_0=\sigma_u^2$, d'où les deux formules. Ensuite $\mathrm{Var}(v\mid y)=\Sigma_0-\lambda^2\mathrm{Var}(y)=\Sigma_0-\Sigma_0/2$,
et le profit espéré vaut $\beta\Sigma_0-\lambda\beta^2\Sigma_0=\sqrt{\Sigma_0}\sigma_u/2$.
\end{proof}

La pente $\lambda$ est le lambda de Kyle (livre~7, chapitre~9)~: l'impact sur le prix d'une unité de flux d'ordres, l'inverse de la profondeur du marché. Davantage de bruit
rend le marché plus profond et l'initié plus riche, et laisse le prix exactement aussi informatif. Sur 100\,000 tirages simulés avec $\sqrt{\Sigma_0}=1$ et
$\sigma_u=2$, la régression de $p$ sur $y$ donne 0.250, la variance résiduelle de $v$ vaut 0.498 et le profit moyen de l'initié 0.991, à comparer aux
valeurs 0.25, 0.5 et 1 de la proposition. Dans la version en temps continu de Kyle, l'initié négocie progressivement, $\lambda$ est constant dans le temps et toute son
information est dans le prix à la fin (\admitted~; démontré dans l'article de Kyle).

\section{Mesurer le trading informé}

\begin{definition}[Probabilité de transaction informée]\label{def:mx:why-there-is-a-spread:pin}
La \emph{probabilité de transaction informée}\index{probabilité de transaction informée} (probability of informed trading, PIN) d'Easley, Kiefer, O'Hara et Paperman est, dans leur modèle
du flux d'ordres quotidien, la part des transactions provenant de \omterm{def:mx:why-there-is-a-spread:informed}{traders informés}~: $\mathrm{PIN}=\alpha\mu/(\alpha\mu+\varepsilon_b+\varepsilon_s)$, où un événement
informationnel survient un jour donné avec probabilité $\alpha$, est une mauvaise nouvelle avec probabilité $\delta$, génère des ordres informés au taux $\mu$, et les achats et ventes non informés arrivent
aux taux $\varepsilon_b$ et $\varepsilon_s$.
\end{definition}

Les achats et les ventes de chaque jour suivent des lois de Poisson dont les taux sont fixés par l'état caché du jour, et les paramètres sont estimés par maximum de vraisemblance sur les comptages quotidiens (la
vraisemblance mêle trois couples de Poisson~; son logarithme doit être calculé de manière numériquement stable).

\omcode{code/firm/spreadmodels/firm_spreadmodels.py}{106}{113}{Log-vraisemblance du PIN, stable grâce au log-sum-exp.}\label{lst:mx:why-there-is-a-spread:pin}

Le modèle interprète comme information toute oscillation commune des achats et des ventes dès lors qu'elle n'est pas unilatérale dans ses données, et Duarte et Young (2009) ont constaté
que l'association du PIN avec les rendements provenait de sa composante d'illiquidité plutôt que de sa composante d'information.
La \Cref{fig:mx:why-there-is-a-spread:pin} montre pourquoi sur des données simulées. Les jours sans information aucune, dont les taux non informés évoluent ensemble d'un jour à l'autre
(un facteur lognormal commun aux deux) sont ajustés avec un PIN de 0.085 lorsque l'écart type du facteur vaut 0.25, de 0.153 à 0.5 et de 0.207 à 0.75, contre 0.006
sans ces oscillations. Un marché de PIN véritable 0.107 est estimé à 0.107 sans oscillations et à 0.195 avec la plus grande.

\begin{omfigure}
\begin{tikzpicture}
\begin{axis}[width=9cm,height=5.2cm,xlabel={\small écart type du facteur d'activité quotidien},ylabel={\small PIN ajusté},
  tick label style={font=\footnotesize},xtick={0,0.25,0.5,0.75},ymin=0,ymax=0.25,grid=major,grid style={gray!25},
  legend columns=2,legend style={font=\footnotesize,at={(0.5,-0.3)},anchor=north,draw=none,fill=none,
  /tikz/every even column/.append style={column sep=8pt}},table/col sep=comma]
\addplot[omDef,very thick,mark=*,error bars/.cd,y dir=both,y explicit] table[x=sd,y=none,y error=none_sd]{figdata/microstructure/04-why-there-is-a-spread/pin.csv};
\addplot[omThm,very thick,mark=square*,error bars/.cd,y dir=both,y explicit] table[x=sd,y=info,y error=info_sd]{figdata/microstructure/04-why-there-is-a-spread/pin.csv};
\addplot[omThm,dashed,domain=0:0.75] {0.1071};
\legend{aucune information (PIN véritable 0),information (PIN véritable 0.107)}
\end{axis}
\end{tikzpicture}

\omcaption{Estimation par maximum de vraisemblance du PIN sur des échantillons simulés de 250 jours (moyenne et écart type sur 20 échantillons), en fonction de la taille d'un facteur quotidien
commun sur les taux non informés~; la ligne en pointillés est le PIN véritable du marché informé. Données~: \texttt{mx\_spread.pin\_bias}.}\label{fig:mx:why-there-is-a-spread:pin}
\end{omfigure}

\section{Tutoriel~: trois fourchettes en code}\label{tut:mx:why-there-is-a-spread:tut}

\begin{tutorial}
\textbf{Objectif.} Calculer les trois fourchettes, vérifier l'équilibre de Kyle par simulation et observer comment le PIN confond activité et information.
\textbf{État final~:} les \cref{fig:mx:why-there-is-a-spread:gm,fig:mx:why-there-is-a-spread:pin} et les nombres des sections 1 et 4.
\begin{tutsteps}
\item \textbf{Rebond.} \texttt{mx\_spread.roll\_on\_tape()}~: l'estimation de Roll, la fourchette affichée (quoted spread) et l'autocorrélation du signe de la transaction sur l'heure simulée.
\item \textbf{Inventaire.} \texttt{inventory\_quotes(v, q, size, gamma, sigma, horizon)} pour des inventaires allant de $-5$ à $5$~: les cotations se déplacent à l'opposé de la
position.
\item \textbf{Apprentissage.} \texttt{gm\_learning()} et \texttt{gm\_with\_cost(mu, dv, c)}.
\item \textbf{Kyle.} \texttt{kyle\_check()} simule l'équilibre et retrouve $\lambda$.
\item \textbf{PIN.} \texttt{pin\_bias(sd, alpha)} sur la grille~; dessiner avec \texttt{fig\_spread.py}.
\end{tutsteps}
\textbf{Pour aller plus loin.} Faire en sorte que le \omterm{def:mx:why-there-is-a-spread:informed}{trader informé} du \omterm{def:mx:why-there-is-a-spread:gm}{modèle de Glosten--Milgrom} ne négocie que lorsque la cotation est assez éloignée de la valeur
(un seuil de profit) et observer la fourchette~; ajuster le PIN sur les jours d'achats et de ventes de \texttt{firm.tape}, où le drapeau informé est connu.
\end{tutorial}

\section{Construction~: les modèles de fourchette}\label{bld:mx:why-there-is-a-spread:spreadmodels}

\begin{build}
\bfield{Rôle} Les implémentations de référence des théories de la fourchette, utilisées par les décompositions du chapitre~5, par le marché à agents du chapitre~27
(un agent informé et un teneur de marché de Glosten--Milgrom) et par les tutoriels sur la tenue de marché du livre~11, comme indices de référence.
\bfield{Interface} \texttt{gm\_quotes}, \texttt{gm\_update}, \texttt{gm\_path}~; \texttt{kyle\_one\_period}, \texttt{kyle\_simulate}~;
\texttt{inventory\_quotes}~; \texttt{roll\_spread}~; \texttt{pin}, \texttt{pin\_loglik}, \texttt{pin\_fit}, \texttt{simulate\_days}.
\bfield{Règles} Formes closes lorsqu'elles existent~; la vraisemblance du PIN sous forme log-sum-exp, avec les log-probabilités de Poisson calculées à partir de \texttt{gammaln}~; ajustements à partir de plusieurs
points de départ~; tout est initialisé à partir d'une graine fixe.
\bfield{Tests d'acceptation} \texttt{code/firm/spreadmodels/tests/}~: la fourchette de Glosten--Milgrom égale à $\mu$ pour $\delta=\tfrac12$ et le cours vendeur égal à la
distribution a posteriori après un achat~; le $\lambda$ et le $\beta$ de Kyle, la variance a posteriori et le profit, ainsi que leurs valeurs simulées~; les cotations d'inventaire calculées à la main~; l'estimateur de Roll sur
une marche aléatoire rebondissante~; le PIN retrouvé à 0.02 près sur 300 jours simulés.
\bfield{Pour aller plus loin} La récursion multi-périodes de Kyle~; le modèle d'Easley--O'Hara à taux d'arrivée variables dans le temps~; le PIN ajusté de Duarte et Young.
\end{build}

\begin{omsources}
\item R.~Roll, ``A simple implicit measure of the effective bid-ask spread in an efficient market'', \emph{Journal of Finance} 39(4), 1984.
\item M.~B.~Garman, ``Market microstructure'', \emph{Journal of Financial Economics} 3(3), 1976.
\item H.~R.~Stoll, ``The supply of dealer services in securities markets'', \emph{Journal of Finance} 33(4), 1978.
\item T.~Ho and H.~R.~Stoll, ``Optimal dealer pricing under transactions and return uncertainty'', \emph{Journal of Financial Economics} 9(1), 1981.
\item L.~R.~Glosten and P.~R.~Milgrom, ``Bid, ask and transaction prices in a specialist market with heterogeneously informed traders'', \emph{Journal of
Financial Economics} 14(1), 1985.
\item A.~S.~Kyle, ``Continuous auctions and insider trading'', \emph{Econometrica} 53(6), 1985.
\item D.~Easley, N.~M.~Kiefer, M.~O'Hara and J.~B.~Paperman, ``Liquidity, information, and infrequently traded stocks'', \emph{Journal of Finance} 51(4), 1996.
\item J.~Duarte and L.~Young, ``Why is PIN priced?'', \emph{Journal of Financial Economics} 91(2), 2009.
\end{omsources}

\section{Exercices}

\begin{exercise}[$\star$]\label{exo:mx:why-there-is-a-spread:1}
Calculer la fourchette que donne l'estimateur de Roll lorsque les variations successives du prix de transaction d'une action ont une autocovariance de $-0.0004$ (en dollars au carré).
\end{exercise}

\begin{exercise}[$\star$]\label{exo:mx:why-there-is-a-spread:2}
Dans le \omterm{def:mx:why-there-is-a-spread:gm}{modèle de Glosten--Milgrom}, la valeur vaut 99 ou 101 avec une probabilité égale, et 20~\% des arrivées sont informées. Calculer le cours acheteur et le cours vendeur.
\end{exercise}

\begin{exercise}[$\star$]\label{exo:mx:why-there-is-a-spread:3}
Dans le \omterm{def:mx:why-there-is-a-spread:kyle}{modèle de Kyle}, $\sqrt{\Sigma_0}=2$ et $\sigma_u=8$. Donner $\lambda$, $\beta$, le profit espéré de l'initié et la variance a posteriori de la valeur.
\end{exercise}

\begin{exercise}[$\star\star$]\label{exo:mx:why-there-is-a-spread:4}
Une dealer valorise sa position à 50, est vendeuse de 3 unités, et a $\gamma=0.2$, $\sigma=1$ et $T=0.5$. Calculer son cours acheteur et son cours vendeur pour une unité. Expliquer pourquoi
les deux cours sont supérieurs à 50.
\end{exercise}

\begin{exercise}[$\star\star$]\label{exo:mx:why-there-is-a-spread:5}
Dans l'exercice~2, un ordre d'achat arrive. Calculer la nouvelle probabilité de la valeur haute pour le fournisseur de liquidité, puis les nouvelles cotations.
\end{exercise}

\begin{exercise}[$\star\star$]\label{exo:mx:why-there-is-a-spread:6}
Expliquer pourquoi doubler le volume des \omterm{def:mx:why-there-is-a-spread:informed}{noise traders} dans le \omterm{def:mx:why-there-is-a-spread:kyle}{modèle de Kyle} ne rend pas le prix moins informatif.
\end{exercise}

\begin{exercise}[$\star\star\star$]\label{exo:mx:why-there-is-a-spread:7}
\emph{Programmation.} Estimer la PIN sur 250 jours simulés avec $\alpha=0.4$ et un facteur d'activité quotidien d'écart type 0.5, puis l'estimer à nouveau après avoir divisé les achats et les ventes de chaque
jour par le total de la journée, puis de les multiplier par le total moyen de l'échantillon. Dire si la normalisation supprime le biais et expliquer pourquoi elle ne le fait que partiellement.
\end{exercise}

\begin{exercise}[$\star\star\star$]\label{exo:mx:why-there-is-a-spread:8}
\emph{Trouver la faille.} «~L'action A a une PIN de 0.25 et l'action B une PIN de 0.12, donc les teneurs de marché de A font face à deux fois plus de trading informé.~»
\end{exercise}

\section{Problème~: Les deux centimes du spécialiste}

\begin{problem}[{Problème du week-end --- Les deux centimes du spécialiste}]\label{pb:mx:why-there-is-a-spread:1}
Un spécialiste affiche une fourchette de deux centimes et veut savoir ce qu'elle rémunère~: le traitement, l'inventaire ou l'information~; et un desk de recherche veut savoir si
la PIN peut le dire.

\medskip\noindent\textbf{Partie I --- Traitement et inventaire.}
\begin{enumerate}
\item Préciser ce que suppose l'estimateur de Roll et ce qu'il donne sur l'heure simulée.
\item Expliquer pourquoi il reste en deçà de la fourchette affichée dans ce cas.
\item Énoncer les cotations du dealer à moyenne--variance et la fourchette qu'elles impliquent.
\item Dire ce qu'il advient des cotations lorsque le dealer est en position acheteuse, et expliquer pourquoi.
\end{enumerate}

\medskip\noindent\textbf{Partie II --- Le \omterm{def:mx:why-there-is-a-spread:gm}{modèle à transactions séquentielles}.}
\begin{enumerate}[resume]
\item Dériver le cours vendeur de Glosten--Milgrom pour $\delta=\tfrac12$.
\item Calculer la fourchette avec 30~\% d'arrivées informées et une valeur valant 0 ou 1.
\item Calculer le nombre de transactions nécessaires à la trajectoire médiane pour que la fourchette tombe au dixième, avec 10~\%, 30~\% et 50~\% d'arrivées informées.
\item Avec un coût de traitement de 0.05 par transaction, calculer la part de la fourchette qui relève de la sélection adverse pour $\mu=0.1$ et $\mu=0.3$.
\end{enumerate}

\medskip\noindent\textbf{Partie III --- Le trader stratégique.}
\begin{enumerate}[resume]
\item Dériver le $\lambda$ de Kyle et le $\beta$.
\item Calculer la part de l'information privée que révèle le prix.
\item Dire ce que donnent 100\,000 tirages simulés pour $\lambda$, la variance a posteriori et le profit de l'initié.
\item Expliquer qui paie le profit de l'initié.
\end{enumerate}

\medskip\noindent\textbf{Partie IV --- PIN.}
\begin{enumerate}[resume]
\item Écrire la PIN en fonction des paramètres du modèle et la calculer pour $\alpha=0.4$, $\mu=60$, $\varepsilon_b=\varepsilon_s=100$.
\item Calculer la PIN estimée lorsqu'il n'y a ni information ni variation d'activité.
\item Calculer la PIN estimée en l'absence d'information, avec des variations quotidiennes d'activité d'écart type 0.25, 0.5 et 0.75.
\item Énoncer le \emph{résultat nommé}~: la part de la fourchette due à la sélection adverse, en fonction de la part informée, et le biais de la PIN en présence de variations communes d'activité.
\item Expliquer pourquoi le modèle interprète les variations communes comme de l'information.
\item Dire quelle conclusion Duarte et Young tirent de la PIN et des rendements.
\item Expliquer comment séparer l'illiquidité de l'information dans le flux d'ordres quotidien.
\item Dire en une phrase lequel des trois coûts une fourchette de deux centimes sur une action liquide rémunère surtout.
\end{enumerate}
\end{problem}

\section{Questions d'entretien}

\begin{interviewq}[$\star$ \iqroles{trader, chercheur}]\label{iq:mx:why-there-is-a-spread:1}
Citez les trois composantes de la fourchette et un élément de preuve empirique pour chacune.
\end{interviewq}

\begin{interviewq}[$\star\star$ \iqroles{chercheur}]\label{iq:mx:why-there-is-a-spread:2}
Dérivez le cours vendeur de Glosten--Milgrom lorsque la valeur vaut 0 ou 1 avec une probabilité égale et qu'une fraction $\mu$ des traders est informée.
\end{interviewq}

\begin{interviewq}[$\star\star$ \iqroles{chercheur}]\label{iq:mx:why-there-is-a-spread:3}
Dans le \omterm{def:mx:why-there-is-a-spread:kyle}{modèle de Kyle}, que devient l'impact sur le prix et le profit de l'initié lorsque le noise trading double~?
\end{interviewq}

\begin{interviewq}[$\star\star$ \iqroles{trader}]\label{iq:mx:why-there-is-a-spread:4}
Vous êtes teneur de marché et vous avez une position acheteuse de 10\,000 actions. Comment déplacez-vous vos cotations, et pourquoi~?
\end{interviewq}

\begin{interviewq}[$\star\star$ \iqroles{chercheur}]\label{iq:mx:why-there-is-a-spread:5}
L'estimateur de Roll donne une fourchette imaginaire pour une action. Que s'est-il passé, et que faites-vous~?
\end{interviewq}

\begin{interviewq}[$\star\star\star$ \iqroles{chercheur}]\label{iq:mx:why-there-is-a-spread:6}
Comment estimeriez-vous la part de la fourchette d'une action qui rémunère la sélection adverse~?
\end{interviewq}
